% GENERATED FILE. Edit canonical lessons, then run npm run generate:books.
% canonical-source-hash: fnv1a64:3e71b5a7170d8093
% canonical-lessons: ZH-C73-xiache, ZH-W73-xia, ZH-C73-shangche, ZH-C73-huochezhan, ZH-C73-jichang, ZH-C73-piao

\chapter{At the Station}
\label{ch:zh-at-the-station}
\hlchaptermodality{73}

\begin{chapteropening}
\textbf{By the end of this chapter:} \emph{I can say get on and get off, and name the station, the airport and tickets.}

Complete the last lesson of chapter 73: I can say get on and get off, and name the station, the airport and tickets.
\end{chapteropening}

\section[\texorpdfstring{\zh{下车} (xiàchē)}{下车 (xiàchē)}]{\zh{下车} (xiàchē) — to get off (a bus or train)}

\label{lesson:ZH-C73-xiache}



\begin{quote}
Before the new one: say the Mandarin for road, then the Mandarin for to arrive.
\end{quote}

\subsection*{You'll want to know: \zh{下车}}
\textbf{\zh{下车}} — \emph{xiàchē} — \textquotedblleft{}to get off (a bus or train)\textquotedblright{}.

\begin{quote}
\textbf{\zh{我们下车吃饭}} \emph{wǒmen xiàchē chī fàn} — Let's get off and eat.
\end{quote}

The new character, in pieces: \zh{下} — a bar, a vertical hanging from it, and a dot to the right: \zh{上} turned over. The next lesson puts it in your hand.

\begin{cousinweb}
\textbf{\zh{下车}} (\emph{xiàchē}) is built from \zh{下} 'to go down' and \zh{车} 'vehicle', both inherited from Middle Chinese, \emph{haeX} and \emph{tsyhae} (Guangyun, in Baxter's transcription).
\end{cousinweb}

\begin{grammarlens}[title={What you've built}]
One more word for times and partings.
\end{grammarlens}

\subsection*{Your turn}
\begin{itemize}
\raggedright
  \item \emph{Say it:} \emph{xiàchē}
  \item \emph{Say it:} \emph{xiàchē}, once more
  \item \emph{Say it:} \emph{wǒmen xiàchē chī fàn} — Let's get off and eat.
  \item \emph{Recall:} say the Mandarin for road, then the Mandarin for to arrive, then say \emph{xiàchē} again
\end{itemize}

\begin{tcolorbox}[breakable,colback=teal!4,colframe=teal!35!black,title={Before you move on}]
What does \zh{下车} mean? (\textquotedblleft{}To get off (a bus or train)\textquotedblright{}.) Say it once more.
\end{tcolorbox}
\section[\texorpdfstring{\zh{下} (xià)}{下 (xià)}]{\zh{下} — a bar, a vertical, a dot}

\label{lesson:ZH-W73-xia}



\begin{quote}
The word first, then the character inside it.

\begin{itemize}
\raggedright
  \item \emph{Recall:} say \emph{xiàchē} — \textbf{R1}, one lesson back, before it decays
\end{itemize}

\begin{quote}
\zh{上} \zh{车}
\end{quote}

\emph{Read it:} each one aloud, with nothing said first — \emph{shàng}, \emph{chē}
\end{quote}

\subsection*{Script — a bar, a vertical, a dot}
\hlblockfigure{\detokenize{figures/ZH-W73-xia-filmstrip.pdf}}{How \zh{下} is written, stroke by stroke}

\begin{quote}
\zh{下}
\end{quote}

\textbf{One piece, three strokes.}

\textbf{A long horizontal} across the top.

\textbf{A vertical} hanging down from its middle.

\textbf{A dot} to the right of the vertical, slanting down to the right.

\subsection*{Your turn}
\begin{itemize}
\raggedright
  \item \emph{Write it:} \zh{下}, all three, with the model in front of you
  \item \emph{Cover:} write it again from memory, and count the pieces, not the strokes
\end{itemize}

\begin{tcolorbox}[breakable,colback=teal!4,colframe=teal!35!black,title={Before you move on}]
How many strokes in \zh{下}? (\textbf{Three}.) Which character you know is it turned over? (\textbf{\zh{上}}.) Where does the dot go? (\textbf{To the right of the vertical}.)

Source: \href{https://raw.githubusercontent.com/chanind/hanzi-writer-data/68d10a4b21150cae5e1ebbd223eed289cf32d90c/data/\%E4\%B8\%8B.json}{Hanzi Writer Data \zh{下}, pinned stroke paths and medians}.
\end{tcolorbox}
\section[\texorpdfstring{\zh{上车} (shàngchē)}{上车 (shàngchē)}]{\zh{上车} (shàngchē) — to get on (a bus or train)}

\label{lesson:ZH-C73-shangche}



\begin{quote}
Before the new one: say the Mandarin for to arrive, then the Mandarin for to get off.
\end{quote}

\subsection*{You'll want to know: \zh{上车}}
\textbf{\zh{上车}} — \emph{shàngchē} — \textquotedblleft{}to get on (a bus or train)\textquotedblright{}.

\begin{quote}
\textbf{\zh{下雨快上车}} \emph{xiàyǔ kuài shàngchē} — It's raining: quick, get in!
\end{quote}

\begin{cousinweb}
\textbf{\zh{上车}} (\emph{shàngchē}) is built from \zh{上} 'to go up' and \zh{车} 'vehicle', both inherited from Middle Chinese, \emph{dzyangX} and \emph{tsyhae} (Guangyun, in Baxter's transcription).
\end{cousinweb}

\begin{grammarlens}[title={What you've built}]
One more word for times and partings.
\end{grammarlens}

\subsection*{Your turn}
\begin{itemize}
\raggedright
  \item \emph{Say it:} \emph{shàngchē}
  \item \emph{Say it:} \emph{shàngchē}, once more
  \item \emph{Say it:} \emph{xiàyǔ kuài shàngchē} — It's raining: quick, get in!
  \item \emph{Recall:} say the Mandarin for to arrive, then the Mandarin for to get off, then say \emph{shàngchē} again
\end{itemize}

\begin{tcolorbox}[breakable,colback=teal!4,colframe=teal!35!black,title={Before you move on}]
What does \zh{上车} mean? (\textquotedblleft{}To get on (a bus or train)\textquotedblright{}.) Say it once more.
\end{tcolorbox}
\section[huǒchēzhàn]{huǒchēzhàn — railway station}

\label{lesson:ZH-C73-huochezhan}



\begin{quote}
Before the new one: say the Mandarin for to get off, then the Mandarin for to get on.
\end{quote}

\subsection*{You'll want to know: huǒchēzhàn}
\emph{huǒchēzhàn} — \textquotedblleft{}railway station\textquotedblright{}.

\begin{quote}
\emph{huǒchēzhàn hěn yuǎn} — The station is a long way off.
\end{quote}

Its characters come later; for now it is a word for the ear.

\begin{cousinweb}
\emph{huǒchēzhàn} is built from \zh{火车} 'train' and zhàn 'a stop', inherited from Middle Chinese \emph{treamH} (Guangyun, in Baxter's transcription), which the dictionary glosses as the everyday word for standing alone. The sense 'station' is not given there, so how it got it is not settled.
\end{cousinweb}

\begin{grammarlens}[title={What you've built}]
One more word for times and partings.
\end{grammarlens}

\subsection*{Your turn}
\begin{itemize}
\raggedright
  \item \emph{Say it:} \emph{huǒchēzhàn}
  \item \emph{Say it:} \emph{huǒchēzhàn}, once more
  \item \emph{Say it:} \emph{huǒchēzhàn hěn yuǎn} — The station is a long way off.
  \item \emph{Recall:} say the Mandarin for to get off, then the Mandarin for to get on, then say \emph{huǒchēzhàn} again
\end{itemize}

\begin{tcolorbox}[breakable,colback=teal!4,colframe=teal!35!black,title={Before you move on}]
What does huǒchēzhàn mean? (\textquotedblleft{}Railway station\textquotedblright{}.) Say it once more.
\end{tcolorbox}
\section[jīchǎng]{jīchǎng — airport}

\label{lesson:ZH-C73-jichang}



\begin{quote}
Before the new one: say the Mandarin for to get on, then the Mandarin for railway station.
\end{quote}

\subsection*{You'll want to know: jīchǎng}
\emph{jīchǎng} — \textquotedblleft{}airport\textquotedblright{}.

\begin{quote}
\emph{wǒ qù jīchǎng} — I'm going to the airport.
\end{quote}

Its characters come later; for now it is a word for the ear.

\begin{cousinweb}
\emph{jīchǎng} is built from \zh{机} 'machine' and chǎng 'open ground', both inherited from Middle Chinese, \emph{kj+j} and \emph{drjang} (Guangyun, in Baxter's transcription); the dictionary glosses the second as a threshing floor. The machine here is the plane.
\end{cousinweb}

\begin{grammarlens}[title={What you've built}]
One more word for times and partings.
\end{grammarlens}

\subsection*{Your turn}
\begin{itemize}
\raggedright
  \item \emph{Say it:} \emph{jīchǎng}
  \item \emph{Say it:} \emph{jīchǎng}, once more
  \item \emph{Say it:} \emph{wǒ qù jīchǎng} — I'm going to the airport.
  \item \emph{Recall:} say the Mandarin for to get on, then the Mandarin for railway station, then say \emph{jīchǎng} again
\end{itemize}

\begin{tcolorbox}[breakable,colback=teal!4,colframe=teal!35!black,title={Before you move on}]
What does jīchǎng mean? (\textquotedblleft{}Airport\textquotedblright{}.) Say it once more.
\end{tcolorbox}
\section[piào]{piào — ticket}

\label{lesson:ZH-C73-piao}



\begin{quote}
Before the new one: say the Mandarin for railway station, then the Mandarin for airport.
\end{quote}

\subsection*{You'll want to know: piào}
\emph{piào} — \textquotedblleft{}ticket\textquotedblright{}.

\begin{quote}
\emph{wǒ mǎi le piào} — I've bought the tickets.
\end{quote}

Its characters come later; for now it is a word for the ear.

\begin{cousinweb}
\emph{piào} is written with a character the Guangyun records, under an older form, as \emph{phjiew} (Guangyun, in Baxter's transcription), glossed 'flying fire'. That is probably a different word: the ticket sense is not given there, and \emph{phjiew} would regularly give piāo, not piào. How the character came to write the word for ticket is not settled.
\end{cousinweb}

\begin{grammarlens}[title={What you've built}]
One more word for times and partings.
\end{grammarlens}

\subsection*{Your turn}
\begin{itemize}
\raggedright
  \item \emph{Say it:} \emph{piào}
  \item \emph{Say it:} \emph{piào}, once more
  \item \emph{Say it:} \emph{wǒ mǎi le piào} — I've bought the tickets.
  \item \emph{Recall:} say the Mandarin for railway station, then the Mandarin for airport, then say \emph{piào} again
\end{itemize}

\begin{tcolorbox}[breakable,colback=teal!4,colframe=teal!35!black,title={Before you move on}]
What does piào mean? (\textquotedblleft{}Ticket\textquotedblright{}.) Say it once more.
\end{tcolorbox}
